\begin{abstract}
Frozen SAM-style segmenters often fail on texture partitions, but failure does not by itself show that texture evidence is absent. We study texture segmentation as a recoverability problem: how much task value can be extracted from frozen features and frozen proposals before retraining the backbone or proposal generator. Our main contribution is \proproute{}, instantiated by \sys{}, which keeps the proposal bank fixed and learns texture partition commitment through proposal compatibility scoring, conservative core selection, and an RWTD-specific dense rescue layer. We retain a smaller \featroute{} branch only as auxiliary diagnostic evidence for latent texture structure in coarse frozen features. Under matched evaluation, \sys{} improves RWTD common-253 coherence from 0.6163 to 0.7013 ARI against a public TextureSAM rerun while keeping mIoU comparable, and improves STLD common-182 from 0.5140/0.7526 to 0.7195/0.7791. Oracle analysis shows that the frozen bank still contains substantial unrecovered value: on RWTD the bank upper bound reaches 0.5183 mIoU / 0.8580 ARI. These results support a narrower claim: in frozen SAM-style texture segmentation, the main missing ingredient is often partition commitment above fragmented proposals rather than new texture features inside the backbone. \codelinksentence
\end{abstract}
